\section{Related Work}
\label{sec:related}

We organise prior work into three threads that BioShield draws on directly:
machine-generated text (MGT) detection, adversarial training for text
classifiers, and LLM-as-agent rewriting.

\subsection{Machine-Generated Text Detection}
Early MGT detectors relied on stylometric and surface-statistical signals.
GLTR \cite{gehrmann2019gltr} visualises token-level rank statistics from a
language model to surface unusually probable continuations. Grover
\cite{zellers2019defending} trained a discriminator jointly with a generative
news LM and showed that the best detector for a given LM is often a model
of the same family — an observation we revisit at biomedical scale.
DetectGPT \cite{mitchell2023detectgpt} and zero-shot likelihood-perturbation
methods exploit the curvature of the source LM's likelihood landscape.
RoBERTa-based supervised detectors \cite{solaiman2019releasestrategies} remain
the most widely deployed baseline, including the GPT-2 output detector and
its biomedical-domain descendants \cite{liang2023gptdetect}. Recent
benchmarks RAID \cite{dugan2024raid} and MAGE \cite{li2024mage} show
that almost all such detectors collapse on out-of-distribution generators,
which our cross-domain RAID probe confirms.

\subsection{Adversarial Training for Text Classifiers}
Adversarial training in NLP spans two distinct families. \emph{Sample-level}
attacks such as TextAttack \cite{morris2020textattack},
TextFooler \cite{jin2020textfooler}, and BERT-Attack
\cite{li2020bertattack} perturb individual tokens or sentences against a
fixed classifier. \emph{Generator-level} adversarial training instead
trains a generator alongside the classifier; SeqGAN \cite{yu2017seqgan}
remains the canonical baseline using policy-gradient REINFORCE. More recent
work on RL-tuned generators \cite{ouyang2022rlhf} and DPO-style preference
optimisation \cite{rafailov2023dpo} suggests that hardening generators via
explicit reward signals is more effective than the SeqGAN regime — which
our flat Condition~B curve indirectly supports.

A practical question that this literature largely sidesteps is
\emph{checkpoint selection}: when retraining yields a sequence of detector
weights with comparable validation metrics, which checkpoint is shipped?
Best-validation-AUC selection is the silent default in most public
implementations. Our Section~\ref{sec:discussion} demonstrates that this
default actively hides retraining signal when the baseline is already at
ceiling, producing the C\,$=$\,D pathology we report.

\subsection{LLM-as-Agent Rewriting and Red-Teaming}
LLMs are increasingly used as \emph{agents} that probe and harden classifiers
by paraphrase or rewrite \cite{perez2022redteaming,ganguli2022redteaming}.
PAIR \cite{chao2024jailbreaking}, AutoDAN \cite{liu2023autodan}, and Tree of
Attacks \cite{mehrotra2024toa} apply this idea to safety classifiers.
Closer to our setting, recent work uses GPT-4 or open-weight LLMs to
paraphrase machine-generated essays until they evade detection
\cite{krishna2023paraphrasing,sadasivan2024aigeneratedtext}. Our agent
loop instantiates the same idea with Qwen2.5-7B as the rewriter and the
top-25\% hardest fakes from the prior round as the agent's input. The
near-zero evasion gain we observe (Condition~C vs. Condition~A) is
consistent with the negative scaling result in
\cite{sadasivan2024aigeneratedtext} that paraphrase-based attacks plateau
when the source generator is already in-family for the detector.

\subsection{Position of BioShield}
BioShield's contribution is not a new architecture or attack. It is a
controlled, fully reproducible \emph{biomedical-domain} measurement of how
much of the evasion-rate signal in published detector papers can be
attributed to (a)~the choice of generator family, (b)~retraining rounds, and
(c)~rewrite agents — and the answer, at \(3\)k scale, is that (a)~dominates
by two orders of magnitude over (b) and (c). This complements the
cross-domain failure mode reported by RAID \cite{dugan2024raid} and MAGE
\cite{li2024mage} with an in-domain explanation for why the same detector
can look near-perfect or near-useless depending on which generator the
reviewer probes it with.
